\subsection{Árbol principal definitivo}
\label{app:prototipo_A7}
Se mide un diámetro exterior 12,7\,mm y longitud 235\,mm. Las tuercas CAD son
de familia 1/2--13, coherentes nominalmente con la varilla de media pulgada;
no se considera el cilindro como evidencia de una varilla lisa ni se
deduce de esa designación el grado o diámetro mínimo resistente real.
Las estaciones medidas en el marco del solver son:
\begin{center}\begin{tabular}{ll}\toprule
Estación & Coordenada axial [mm]\\\midrule
Guía de placa 16 & -94,6302\\
Centro de agujero del rodamiento 53 & -153,4231\\
Unión manivela--árbol & -245,7550\\
Centro de empuñadura & -307,7550\\
\bottomrule\end{tabular}\end{center}

La separación entre guías es 58,793\,mm.
La fuerza de mano se traslada a su unión real como fuerza y par flector,
manteniendo la excentricidad axial 62\,mm. No se prolonga ficticiamente el
árbol hasta el centro de la mano. Las ecuaciones nominales son:
\begin{align}
 \sigma_N&=4F/(\pi d^2),&\tau_T&=16T/(\pi d^3),\\
 \sigma_{VM}&=\sqrt{\sigma_N^2+3\tau_T^2},&
 \delta_N&=FL/(EA),\quad\theta=TL/(GJ).
\end{align}
La flexión se calcula por la vigaEB de EP2 con las dos guías radiales
ideales, giros libres, y cargas/par de manivela en su estación efectiva.
Se mantiene $E=210$ GPa y $\nu=0,30$ de referencia, sin atribuirles una clase
resistente. La envolvente flectora combina orientaciones posibles dentro
de un mismo estado correlacionado, sin mezclar máximos de otros ángulos.
\begin{center}\begin{tabular}{ll}\toprule
Magnitud nominal & Resultado\\\midrule
Compresión máxima & 7,247 MPa\\
Cortante de torsión máximo & 6,912 MPa\\
Von Mises axial--torsión, correlacionado & 13,994 MPa\\
Envolvente con flexión & 23,460 MPa\\
Acortamiento sobre 235\,mm & 0,00811 mm\\
Giro sobre 235\,mm & 0,1815 grados\\
Flecha transversal de referencia & 0,06106 mm\\
Reacción radial manual, apoyo inferior & 61,076 N\\
\bottomrule\end{tabular}\end{center}

La sección de 12,7\,mm es la referencia nominal exterior: faltan la raíz
efectiva y el grado de la varilla para dar un factor de seguridad real,
igual que en la revisión precedente. No se traslada la clase 8.8 del M8
a este árbol. Si todo el empuje se transmite por él, el apoyo debe recibir
924,32\,N, con pesos de referencia.
El buje transmite carga axial solamente si el árbol queda retenido en él;
una coincidencia cilíndrica 12,7/12,7 no demuestra esa retención.
